Hybrid quantum--classical models are increasingly proposed for financial
prediction, but the reported gains are rarely separated from the classical
components they are wrapped in, or are measured on windows too short to
resolve them. We study a hybrid quantum temporal convolutional
network for daily cross-sectional ranking of \num{2000} Tokyo Stock Exchange
equities under a protocol built to make the quantum contribution identifiable:
five conditions (trained circuit, frozen circuit, parameter-matched classical
surrogate, no circuit, circuit only) trained through one code path; four
rolling-origin walk-forward folds with an embargo and five seeds; annualised
metrics with stationary-bootstrap intervals, Deflated Sharpe Ratios,
backtest-overfitting probabilities and family-wise error control; five
classical baselines including a gradient-boosted tree; and a locked holdout
evaluated once. On \num{500} out-of-sample days two one-sided tests place the
trained-versus-frozen difference inside \num{0.03} annualised Sharpe
($p < \num{0.001}$, minimum detectable effect \num{0.02}): the circuit is demonstrably equivalent to a frozen one, not merely
indistinguishable from it. Against the network with no circuit the design
resolves only a \num{0.50} margin, and we say so. No condition is
distinguishable from zero, the best Deflated Sharpe Ratio is \num{0.012}, and
nine of the ten have no break-even transaction cost because their gross return
is negative. A locked holdout on a profitable window reproduces the ablation, and so do
\num{2000} US equities over the same calendar, where every condition loses
money and the same equivalence holds ($p = \num{0.0005}$). The circuit is
trainable, with no barren plateau to twelve qubits and rankings robust to gate
noise, and a second ansatz family with data re-uploading and a global
observable gives the same ordering, its difference bounded at \num{0.25}
Sharpe. Measuring the shot budget at the input distribution the encoder
actually produces raises the cost of preserving the ranking by two orders of
magnitude, to about \num{32768} shots per security. Re-scoring our own previously published result on this dataset over
\num{241} out-of-sample days rather than the \num{22} it reported leaves it
ranked first but no longer significant, and traces its advantage to the
window, single-seed selection, a reversed split adjustment and the classical
encoder. We
release the batched circuit implementation and the calibrated evaluation
stack.
